%!TEX root = ./../../springer-mono-format/book.tex
%!TEX root = ./../../thesis-trier.tex

\paragraph{Partial Functions.} 
In principle, everything that can be formulated with partial functions
can be recast using only (total) functions, and every PDFA
can be converted into an automaton as defined before by adding a trap
state at which all previously undefined transitions end. 
The reader might therefore ask what is gained from this generalization. 
In theoretical computer science, especially in computability theory (see~\cite{HopUll79,DBLP:books/daglib/0088160}), where a notion of effectiveness (or computability) is associated with functions via input-output machines, it is often natural to allow the ``result'' to be undefined. In this setting, a function may be undefined for inputs on which a machine is not required to produce a result, or on which it may fail to terminate by entering an infinite loop.
In our context, we introduced PDFAs based on partial functions.  
As we will also see in the next paragraph, it is convenient to specify specific instances of the constrained
synchronization problem by PDFAs.

\paragraph{Partial Automata.} We will not be concerned with (carefully or exactly)
synchronizing PDFAs in this work. The reason for introducing partial automata
is that they give a convenient specification for constraint languages.
More specifically:
\begin{enumerate}
	\item If $\mathcal A$ is a complete automaton having only coaccessible states, 
	then $\constrsync{L(\mathcal A)}$
	is equivalent to~$\constrsync{\Sigma^*}$
	see~\cref{constrsync:sec:constraut:unconstrained}, 
	\crefitem{constrsync:cor:upperset}{constrsync:cor:upperset:item:complete}, and \cref{rem:PDFA-for-constraint} for details.
	But non-coaccessible states can be discarded to give
	a PDFA accepting the same language.
    So, instead of carrying around these useless states it is more economical 
	to use trim PDFAs, in particular in our discussions of small constraint automata later in \cref{constr_sync:sec:two-state,sec:constr_sync:three_states}. 
	
	\item There is some freedom on how to distribute the start state and the final states 
	without altering the computational complexity of the corresponding constrained synchronization problem (see \cref{constr_sync:prop:final_states}). 
	This might give rise to many different regular expressions 
	for the resulting languages while the underlying automaton, apart from its start state and set of final states,
	remains the same.
	Because of this, we regard PDFAs as a better specification mechanism compared to regular expressions.
\end{enumerate}

The fact mentioned in \cref{constrsync:rem:alphabet}
that we can ignore letters that are not used in any transitions up to $\AC^0$-reductions
was stated \cite[Lem.~11]{FernauGHHVW19} for $\lelogm$-reductions.


\IN{synchronizing prefix code}
\paragraph{Codes and Overlap-Free Languages.}
Codes, and a connection to synchronizing automata,
appeared in our introductory
example on error-prone communication from \cref{intro:sec:examples:communication}.
In this example, a decoder modeled by a synchronizing automaton
allows recovery from errors by sending a synchronizing word
that realigns the decoder with the encoder.
This connection to codes has been noted since the early days
of synchronizing automata,
for example, by Liu~\cite{Liu63} and Laemmel~\cite{Laemmel56,Laemmel63,Laemmel69};
see also the surveys~\cite{KariVolkov2021,Volkov22},
the textbook~\cite{DBLP:books/daglib/0025093} 
or the theses~\cite{Biskup2008,DBLP:phd/hal/Ryzhikov20}.

Without reference to automata or regular languages,
Schützenberger~\cite{DBLP:journals/iandc/Schutzenberger67} called a prefix code $C \subseteq \Sigma^+$
\emph{synchronizing} if there exists $u \in C^*$
such that $\Sigma^* u \subseteq C^*$.\footnote{This definition implies 
	maximality of the code~$C$ w.r.t.~set inclusion.
	For consider $v \in \Sigma^*$.
	By assumption, $vu \in C^*$.
	Thus this word has a factorization
	in words from~$C$, which excludes
	that $C \cup \{v\}$ can be a code
	since $u \in C^*$ implies that it also 
	has a factorization in words from $C \cup \{v\}$.
	For a more general definition of synchronization applying to a wider class of prefix codes (and even general codes), see~\cite{DBLP:books/daglib/0025093}.}
Almost every prefix code that is maximal for set-inclusion
is synchronizing~\cite{DBLP:journals/tit/FreilingJTZ03}. 

We can obtain a decoding automaton for a regular prefix code $C$
via the constructions outlined in
\cref{syncmax:rem:root-tree-prefix-code,constr_sync:prop:aut_prefix_free_exact_sync,syncmax:rem:general-aut-prefix-code}.
Each time the decoding automaton enters a final state, this indicates
the detection of a code word in the input stream.
According to \cref{constr_sync:prop:aut_prefix_free_exact_sync},
this automaton admits exactly synchronizing words, and each such word
can be extended to an exactly synchronizing word ending in the unique
final state.
In case the automaton is complete, each exactly synchronizing word
is a synchronizing word in the usual sense as defined in \cref{part:intro},
Moreover, those exactly synchronizing words ending in the final state
are precisely the synchronizing words of~$C$.
For example, consider the complete automaton obtained from
\cref{intro:subfig:decoder}, whose start state and unique final state
coincide at $q_{\varepsilon}$.
The word $0110$, which synchronizes every state to $q_{\varepsilon}$,
is also a synchronizing word of the code.
Apart from the literature already mentioned above, for more information
on the relation of codes to synchronizing automata,
see also~\cite{BeaBerPer2011,DBLP:journals/tcs/CarpiD17,DBLP:conf/dlt/GusevP11,DBLP:phd/hal/Ryzhikov20}.

Overlap-free languages were introduced in~\cite{DBLP:conf/cocoon/HanW06},
but with a different definition
by defining~$L \subseteq \Sigma^*$ overlap-free if for all distinct $u,v \in L$
we have that $\pref{u} \cap \suff{v}$ and $\suff{u} \cap \pref{v}$
both equal $\{\varepsilon\}$ and $\pref{u} \cap \suff{u} = \{\varepsilon, u\}$
for all $u \in \Sigma^*$.\footnote{In fact, the original definition
demands equality to the empty set of $\pref{u}\cap \suff{v}$, which does not seem to be the intended
definition since the empty word is always a common prefix and suffix
of any two words,  and also the condition
$\pref{u}\cap \suff{u} = \{\varepsilon, u\}$ is missing that excludes a non-empty proper overlap
of a word with itself only a condition on distinct words is imposed. See also the footnote in~\cite[p.~33]{DBLP:journals/tcs/JurgensenS21} which
remarks on the problematic definition from~\cite{DBLP:conf/cocoon/HanW06}.}
This implies overlap-freeness in our sense, but also that~$L$
is a bifix code since if, without loss of generality,
$v \in \pref{u}$ for distinct $u,v \in \Sigma^*$, then $v \in \pref{u} \cap \suff{v}$ since $v \in \suff{v}$.
For example, $L = \{aba, b\}$ fulfills this definition of overlap-freeness,
but $\{a,ab\}$, which is overlap-free according to our definition,
is excluded with this definition. In terms of \cref{constrsync:def:codes}
this wider definition corresponds to replacing the condition
``$u \ne \varepsilon$ and $w \ne \varepsilon$'' from our definition
by ``$u \ne \varepsilon$ or $w \ne \varepsilon$''. We have not seen
this definition of overlap-freeness elsewhere and the one we adopted
seems to be the standard definition used, for example, 
also in~\cite{DBLP:journals/ijfcs/HanS11,DBLP:journals/iandc/Jurgensen08,DBLP:reference/hfl/JurgensenK97,DBLP:journals/tcs/JurgensenS21,DBLP:journals/jalc/KonstantinidisY10}.
In fact, the authors of~\cite{DBLP:conf/cocoon/HanW06}, a conference contribution,
later adopted the definition we gave here in their extended journal version~\cite{DBLP:journals/ijfcs/HanS11}
of the conference contribution.

As noted in \cref{constrsync:sec:codes:solid}, the characterization of solid codes
from \crefitem{thm:solidcode}{thm:solidcode:2} appears to be new.

We note that for codes that are regular languages, often interesting relations hold among the classes
considered in the present work. For example, in~\cite{DBLP:journals/jalc/Konstantinidis20}
it was proved that a regular bifix code is a finite union of regular infix codes.



% dazu noch was sagen? auch noch zu solid codes mehr sagen?
%\paragraph{The Levenshtein-Romanov Mapping.}


%(in particular, we can regard $C$ as a set of abstract symbols and take $\Gamma = C$).
\IN{free submonoid (of $\Sigma^*$)}
\IN{defect theorem}
\paragraph{Monoid Languages.} The free monoid $\Sigma^*$ has a central role in theoretical computer science
as our primary means of representing and processing information. As $\Sigma^*$ is a monoid, languages
that are submonoids of~$\Sigma^*$ form a natural object of study. A submonoid
having a generating set $C$ that is a code is said to be a \emph{free \mbox{(sub-)}monoid} of $\Sigma^*$. 
By \cref{constr_sync:prop:coding_morphisms} a bijection between some set $\Gamma$
and~$C$ extends to a coding morphism.
%\footnote{In particular,
%	for $X = C$ we obtain an isomorphism between $C^*$ considered as 
%	the set of finite sequences	and the monoid generated by $C$ in $\Sigma^*$, and it is isomorphic to 
%	a free monoid over $X$ as defined in \SEC~\ref{sec:formal_language_theory}.}	 
Conversely, the image of a coding morphism
yields a submonoid generated by a code. So, free submonoids of $\Sigma^*$
are precisely the images of coding morphisms, i.e., isomorphic
images of free monoids over some set~$\Gamma$ as
introduced in \cref{sec:formal_language_theory} as sets of finite sequences,
which explains why we call them free monoids as well.
The relevance of these submonoids lies in the transmission, translation and encoding of information w.r.t.~different representations.
Every language over $\Sigma$
is contained in a smallest free submonoid of $\Sigma^*$~\cite{lothaire_1997}.
The \emph{defect theorem}~\cite[Thm. 1.2.5]{lothaire_1997}
states that in case $L \subseteq \Sigma^*$ is a finite language that is not a code,
there exists a finite code $C \subseteq \Sigma^+$
with $|C| \le |L| - 1$ such that $L \subseteq C^*$.

\IN{thin language}
\IN{slender language}
\IN{dense language}

 
\paragraph{The Growth Of Languages and Bounded Languages.}

Berman and Hartmanis~\cite{DBLP:journals/siamcomp/BermanH77}
introduced the notion of sparse languages to study encodings of \NP-complete problems.
They conjectured that, under the assumption $\PTIME \ne \NP$,
sparse languages cannot encode \NP-complete problems with respect to $\lePm$-reductions.
This conjecture was later proved by Mahaney~\cite{DBLP:journals/jcss/Mahaney82}
(see also the survey~\cite{DBLP:conf/mfcs/HartmanisM80}).
Surveys of sparse regular languages can be found
in~\cite[Chap.~2, Sect.~3.5]{DBLP:reference/hfl/Yu97}
and~\cite[Sect.~XII.4.2]{Pin2022}.

P\u{a}un and Salomaa~\cite{DBLP:journals/dam/PaunS95}
introduced \emph{thin} and \emph{slender}
languages where $|L \cap \Sigma^n| \le 1$ and $|L \cap \Sigma^n| \le c$, respectively,
for sufficiently large $n$.
Slender context-free languages were investigated in~\cite{DBLP:journals/tcs/Ilie94,DBLP:journals/tcs/Raz97}.
Bertoni et al.~\cite{DBLP:conf/stacs/BertoniGS87} investigated the computational complexity
of the function $n \mapsto |L \cap \Sigma^n|$
on context-free languages.
D'Alessandro et al.~\cite{DBLP:journals/tcs/DAlessandroIV06} studied this function for sparse context-free languages.

Ginsburg and Spanier~\cite{Ginsburg1964BOUNDEDAL}
introduced bounded languages and investigated bounded context-free languages.
They later studied bounded regular languages
in~\cite{GinsburgSpanier66}.
Several authors extended \cref{constrsync:thm:reg-sparse-or-exp}
to context-free languages~\cite{DBLP:journals/jcss/BridsonG02,DBLP:conf/stacs/IbarraR86,DBLP:journals/tcs/Incitti01,DBLP:journals/eik/LatteuxT84,DBLP:journals/tcs/Raz97,Trofimov81}.
Indexed languages do not satisfy this dichotomy: there
exist indexed languages of intermediate growth~\cite{DBLP:journals/tcs/GrigorchukM99}.
Ginsburg and Spanier~\cite[Thm.~5.2]{Ginsburg1964BOUNDEDAL}
proved that boundedness is decidable for context-free languages.
Consequently, the growth type of a context-free language is decidable.
Gawrychowski et al.~\cite{DBLP:journals/ijfcs/GawrychowskiKRS10}
gave a polynomial-time algorithm for deciding the growth type of a context-free language.

Ilie et al.~\cite{DBLP:journals/ita/IlieRS00}
characterized sparse context-free languages via special Dyck languages.
Szilard et al.~\cite[Cor.~1]{DBLP:conf/mfcs/SzilardYZS92}
characterized sparse regular languages by regular expressions of star-height one.
Shur~\cite{DBLP:conf/csr/Shur08}
related the growth of a regular language to structural properties of its accepting automata.

A language $L \subseteq \Sigma^*$ is called \emph{dense}
if $\Sigma^* w \Sigma^* \cap L \ne \emptyset$
for every $w \in \Sigma^*$~\cite[Sect.~XII.4.1]{Pin2022}.
Regular dense languages over alphabets of size at least two
are never sparse~\cite[Cor.~4.28]{Pin2022}.

Our terminology for letter-bounded languages is not standard.
For example,~\cite{DBLP:journals/iandc/IbarraR16}
uses the term letter-bounded for the strict variant in which all letters are distinct,
whereas~\cite{DBLP:journals/tcs/CarpiDIM21}
refers to this variant as distinct-letter-bounded.

The function~$\rho$ via equal-letter runs occurs naturally in text compression
algorithms, see~\cite{DBLP:conf/dlt/FrosiniMRRS22,Prezza16},
and combinatorics on words~\cite{DBLP:journals/tcs/BiagiCLR25}
since it is related to run-length encodings. In the theory of synchronizing automata,
number of equal-letter runs in synchronizing words have also 
been investigated in~\cite{DBLP:journals/jalc/DonZ19}.